% Part: sets-functions-relations
% Chapter: sets
% Section: pairs-and-products

\documentclass[../../../include/open-logic-section]{subfiles}

\begin{document}

\olfileid{sfr}{set}{pai}
\olsection{Paren, n-tallen en cartesische producten}

\begin{explain}
Uit extensionaliteit volgt dat de elementen van een verzameling geen
volgorde hebben. Om een volgorde weer te geven gebruiken we daarom
\emph{geordende paren} $\tuple{x, y}$. In een ongeordend paar $\{x,
y\}$ doet de volgorde er niet toe: $\{x, y\} = \{y, x\}$. In een
geordend paar is de volgorde wel van belang: als $x \neq y$, dan
$\tuple{x, y} \neq \tuple{y, x}$.

Hoe moeten we geordende paren in de verzamelingenleer opvatten?
Van wezenlijk belang is dat twee geordende paren dan en slechts dan gelijk
zijn als hun eerste elementen gelijk zijn en ook hun tweede elementen
gelijk zijn, dus:
\[
  \tuple{a, b}= \tuple{c, d}\text{ dan en slechts dan als zowel }a = c \text{ als }b=d.
\]
We kunnen geordende paren in de verzamelingenleer definiëren met de
definitie van Wiener--Kuratowski.
\end{explain}

\begin{defn}[Geordend paar]\ollabel{wienerkuratowski}
	$\tuple{a, b} = \{\{a\}, \{a, b\}\}$.
\end{defn}

\begin{prob}
	Bewijs met \olref[sfr][set][pai]{wienerkuratowski} dat $\tuple{a,
	b}= \tuple{c, d}$ dan en slechts dan als zowel $a = c$ als $b=d$.
\end{prob}

\begin{explain}
Nu we een definitie van een geordend paar hebben vastgelegd, kunnen we
daarmee andere verzamelingen definiëren. Soms hebben we bijvoorbeeld
ook geordende rijtjes van meer dan twee objecten nodig, zoals
\emph{drietallen} $\tuple{x, y, z}$, \emph{viertallen} $\tuple{x, y,
z, u}$, enzovoort. We kunnen drietallen opvatten als bijzondere
geordende paren waarvan het eerste element zelf een geordend paar is:
$\tuple{x, y, z}$ is $\tuple{\tuple{x, y},z}$. Hetzelfde geldt voor
viertallen: $\tuple{x,y,z,u}$ is
$\tuple{\tuple{\tuple{x,y},z},u}$, enzovoort. In het algemeen spreken
we van \emph{geordende $n$-tallen} $\tuple{x_1, \dots, x_n}$.

Bepaalde verzamelingen van geordende paren of andere geordende
$n$-tallen zullen van pas komen.
\end{explain}

\begin{defn}[Cartesisch product]
Gegeven verzamelingen $A$ en $B$ is hun \emph{cartesisch product} $A
\times B$ gedefinieerd door
\[
  A \times B = \Setabs{\tuple{x, y}}{x \in A \text{ en } y \in B}.
\]
\end{defn}

\begin{ex}
Als $A = \{0, 1\}$ en $B = \{1, a, b\}$, dan is hun product
\[
A \times B = \{ \tuple{0, 1}, \tuple{0, a}, \tuple{0, b},
    \tuple{1, 1}, \tuple{1, a}, \tuple{1, b} \}.
\]
\end{ex}

\begin{ex}
Als $A$ een verzameling is, schrijven we het product van $A$ met
zichzelf, $A \times A$, ook als~$A^2$. Dit is de verzameling van
\emph{alle} paren $\tuple{x, y}$ waarvoor $x, y \in A$. De
verzameling van alle drietallen $\tuple{x, y, z}$ is $A^3$,
enzovoort. We kunnen dit recursief definiëren:
\begin{align*}
  A^1 & = A\\
  A^{k+1} & = A^k \times A
\end{align*}
\end{ex}

\begin{prob}
Geef alle !!{element}s van $\{1, 2, 3\}^3$.
\end{prob}

\begin{prop}\ollabel{cardnmprod}
Als $A$ uit $n$ !!{element}s bestaat en $B$ uit $m$ !!{element}s, dan
heeft $A \times B$ precies $n\cdot m$ elementen.
\end{prop}

\begin{proof}
Voor elk !!{element}~$x$ van~$A$ zijn er $m$ !!{element}s van de vorm
$\tuple{x, y} \in A \times B$. Laat $B_x = \Setabs{\tuple{x, y}}{y
  \in B}$. Als $x_1 \neq x_2$, verschilt elk paar uit de ene zo
gevormde verzameling van elk paar uit de andere; in het bijzonder geldt $\tuple{x_1, y} \neq
\tuple{x_2, y}$. Daarom is $B_{x_1} \cap B_{x_2} = \emptyset$. Als
$A = \{x_1, \dots, x_n\}$, dan is bovendien $A \times B = B_{x_1}
\cup \dots \cup B_{x_n}$ en heeft dit product dus $n\cdot m$
!!{element}s.

We kunnen dit weergeven door de !!{element}s van~$A \times B$ in een
rooster te zetten:
\[
\begin{array}{rcccc}
  B_{x_1} = & \{\tuple{x_1, y_1} & \tuple{x_1, y_2} & \dots & \tuple{x_1, y_m}\}\\
  B_{x_2} = & \{\tuple{x_2, y_1} & \tuple{x_2, y_2} & \dots & \tuple{x_2, y_m}\}\\
  \vdots & & \vdots\\
  B_{x_n} = & \{\tuple{x_n, y_1} & \tuple{x_n, y_2} & \dots & \tuple{x_n, y_m}\}
\end{array}
\]
Omdat alle $x_i$ onderling verschillend zijn en ook alle $y_j$
onderling verschillend zijn, zijn geen twee paren in dit rooster
gelijk. Er staan er $n\cdot m$ in.
\end{proof}

\begin{prob}
Bewijs met inductie naar~$k$ dat voor alle $k \ge 1$ geldt: als $A$
$n$ !!{element}s heeft, dan heeft $A^k$ precies $n^k$ !!{element}s.
\end{prob}

\begin{ex}
Als $A$ een verzameling is, is een \emph{woord} over~$A$ een
willekeurig rijtje !!{element}s van~$A$. Een rijtje kan worden opgevat
als een geordend $n$-tal van !!{element}s van~$A$. Als bijvoorbeeld $A =
\{a, b, c\}$, dan kan het rijtje ‘$bac$’ worden opgevat als het
drietal~$\tuple{b, a, c}$. Woorden, dat wil zeggen rijtjes symbolen,
zijn van wezenlijk belang in de informatica. Volgens afspraak
beschouwen we !!{element}s van~$A$ als rijtjes van lengte~$1$, en
$\emptyset$ als het rijtje van lengte~$0$. De verzameling van
\emph{alle} woorden over~$A$ is dan
\[
A^* = \{\emptyset\} \cup A \cup A^2 \cup A^3 \cup \dots
\]
\end{ex}
\end{document}
